\chapter[CAPÍTULO V]{COSTO DE IMPLEMENTACIÓN DE MEJORA / INNOVACIÓN / CREATIVIDAD}

\section{Costo de Materiales}

Este rubro reúne los servicios digitales, infraestructura en la nube y herramientas indispensables considerados para poner en funcionamiento y asegurar el acceso público al sistema HelpDesk. A continuación se identifican los servicios registrados, su nombre y sus respectivos costos en soles. Para facilitar la comparación, los importes en dólares se expresan utilizando el tipo de cambio referencial de S/ 3.42 por dólar estadounidense.

\vspace{0.3cm}

\begingroup
\small
\renewcommand{\arraystretch}{1.25}
\begin{longtable}{@{} p{4.5cm} p{2.2cm} p{3.8cm} p{2.2cm} p{2.2cm} @{}}
\caption{{Costo de materiales y servicios digitales}} \label{tab:costo_materiales} \\
\toprule
\textbf{Concepto / Servicio} & \textbf{Nombre} & \textbf{Costo (S/.)} & \textbf{Dólar \$} \\
\midrule
\endfirsthead

\multicolumn{5}{@{}l}{\textbf{Tabla \ref{tab:costo_materiales}} \textit{(Continuación)}} \\[0.4em]
\toprule
\textbf{Concepto / Servicio}  & \textbf{Nombre} & \textbf{Costo (S/.)} & \textbf{Dólar \$} \\
\midrule
\endhead

\midrule
\multicolumn{5}{r}{\textit{Continúa en la siguiente página\dots}} \\
\endfoot

\bottomrule
\multicolumn{5}{@{}p{\textwidth}}{\small \textit{Nota.} Detalle de los servicios digitales y de infraestructura requeridos para la puesta en marcha y despliegue del sistema web. Los importes en dólares se calcularon a partir de los costos en soles, utilizando un tipo de cambio referencial de S/ 3.42 por dólar estadounidense.} \\
\endlastfoot

Base de datos & MySQL / MongoDB / Prisma & S/ 00.00 & \$ 00.00 \\
Servidor Hosting & Servidor Interno J\&P & S/ 0.00 & \$ 0.00 \\
Servidor Local & Memoria Ram Fury Beast 16 GB, 3200 MHz  & S/ 455 & \$ 131.40 \\
Certificado SSL (HTTPS) & PositiveSSL & S/ 88.00 & \$ 25.41 \\
\textbf{Total} & & \textbf{S/ 543.00} & \textbf{\$ 156.81} \\
\end{longtable}
\endgroup

Los servicios digitales considerados comprenden la base de datos en la nube (MySQL / MongoDB / Prisma), el servidor de alojamiento (\textit{hosting}) para el despliegue del backend y frontend, y el registro del dominio web institucional, garantizando la disponibilidad y seguridad de la plataforma para los usuarios y técnicos de J\&P Periféricos S.A.C.

\section{Costo de Mano de Obra}

Los servicios digitales considerados comprenden la base de datos en la nube (MySQL / MongoDB / Prisma), el servidor de alojamiento (\textit{hosting}) para el despliegue del backend y frontend, y el registro del dominio web institucional, garantizando la disponibilidad y seguridad de la plataforma para los usuarios y técnicos de J\&P Periféricos S.A.C. Los importes en dólares se expresan utilizando el tipo de cambio referencial de S/ 3.42 por dólar estadounidense.
\vspace{0.3cm}

\begingroup
\small
\renewcommand{\arraystretch}{1.25}
\begin{longtable}{@{} p{3.9cm} p{1.6cm} p{2.2cm} p{2.5cm} p{2.2cm} p{1.8cm} @{}}
\caption{{Costo de mano de obra del equipo de desarrollo}} \label{tab:costo_mano_obra} \\
\toprule
\textbf{Actividad} & \textbf{Cantidad} & \textbf{Días laborables} & \textbf{Costo unitario (S/.)} & \textbf{Total (S/.)} & \textbf{Dólar \$} \\
\midrule
\endfirsthead

\multicolumn{6}{@{}l}{\textbf{Tabla \ref{tab:costo_mano_obra}} \textit{(Continuación)}} \\[0.4em]
\toprule
\textbf{Actividad} & \textbf{Cantidad} & \textbf{Días laborables} & \textbf{Costo unitario (S/.)} & \textbf{Total (S/.)} & \textbf{Dólar \$} \\
\midrule
\endhead

\midrule
\multicolumn{6}{r}{\textit{Continúa en la siguiente página\dots}} \\
\endfoot

\bottomrule
\multicolumn{6}{@{}p{\textwidth}}{\small \textit{Nota.} Distribución de jornadas (149 días correspondientes a los Sprints 0--5) y costos de mano de obra del equipo de desarrollo. La conversión a dólares se realizó con el tipo de cambio referencial de S/ 3.42.} \\
\endlastfoot

Desarrollador & 1 & 149 & S/ 30.00 & S/ 4,470.00 & \$ 1,307.02 \\
Desarrollador & 1 & 149 & S/ 30.00 & S/ 4,470.00 & \$ 1,307.02 \\
Desarrollador & 1 & 149 & S/ 30.00 & S/ 4,470.00 & \$ 1,307.02 \\
\textbf{Total} & & & & \textbf{S/ 13,410.00} & \textbf{\$ 3,921.05} \\
\end{longtable}
\endgroup
La tarifa diaria de S/ 30.00 por mano de obra considera el perfil junior de los desarrolladores según la clasificación del MTPE(Ministerio de Trabajo y Promoción del Empleo), ajustada a los rangos de mercado para desarrolladores en Lima, Perú, y refleja una tarifa colaborativa institucional por debajo del valor de mercado para fomentar la innovación en J\&P Periféricos S.A.C.; Además, según el Anuario Estadístico MTPE 2024, la remuneración en el sector formal privado promedia S/ 2,898 mensuales, y para jóvenes (18-29 años) alcanza S/ 1,905, lo que respalda que la tarifa de S/ 30.00/día representa un costo racional y ajustado para un perfil junior de desarrollo de software.
\newpage

\section{Costo de Máquinas, Herramientas y Equipos}
\subsection{El Costo de Depreciación de herramientas}

Para desarrollar el sistema HelpDesk, los integrantes del equipo utilizaron sus propias computadoras. Se utiliza la formula: Valor inicial / vida útil en años para calcular el valor residual.

\vspace{0.3cm}

\begingroup
\small
\renewcommand{\arraystretch}{1.25}
\begin{longtable}{@{} p{3.5cm} p{10.9cm} @{}}
\caption{{Equipo de cómputo del primer desarrollador}} \label{tab:equipo_dev1} \\
\toprule
\textbf{Nombre} & \textbf{Descripción} \\
\midrule
\endfirsthead

\multicolumn{2}{@{}l}{\textbf{Tabla \ref{tab:equipo_dev1}} \textit{(Continuación)}} \\[0.4em]
\toprule
\textbf{Nombre} & \textbf{Descripción} \\
\midrule
\endhead

\midrule
\multicolumn{2}{r}{\textit{Continúa en la siguiente página\dots}} \\
\endfoot

\bottomrule
\multicolumn{2}{@{}p{\textwidth}}{\small \textit{Nota.} Computadora utilizada como equipo principal de trabajo por uno de los desarrolladores.} \\
\endlastfoot

Característica de la PC: & AMD Ryzen 5 8600G w/ Radeon 760M Graphics (4.35 GHz) - 32.0 GB RAM \\
Valor inicial: & S/ 2,800.00 \\
Vida útil en años: & 5 años \\
Valor residual: & S/ 560 \\
\end{longtable}
\endgroup

\vspace{0.3cm}

\begingroup
\small
\renewcommand{\arraystretch}{1.25}
\begin{longtable}{@{} p{3.5cm} p{10.9cm} @{}}
\caption{{Equipo de cómputo del segundo desarrollador}} \label{tab:equipo_dev2} \\
\toprule
\textbf{Nombre} & \textbf{Descripción} \\
\midrule
\endfirsthead

\multicolumn{2}{@{}l}{\textbf{Tabla \ref{tab:equipo_dev2}} \textit{(Continuación)}} \\[0.4em]
\toprule
\textbf{Nombre} & \textbf{Descripción} \\
\midrule
\endhead

\midrule
\multicolumn{2}{r}{\textit{Continúa en la siguiente página\dots}} \\
\endfoot

\bottomrule
\multicolumn{2}{@{}p{\textwidth}}{\small \textit{Nota.} Computadora utilizada como equipo principal de trabajo por uno de los desarrolladores.} \\
\endlastfoot

Característica de la PC: & AMD Ryzen 3 2200G - 8.0 GB RAM - NVIDIA GeForce 210 - 224 GB SSD \\
Valor inicial: & S/ 1,600.00 \\
Vida Util en años: & 5 años \\
Valor residual: & S/ 320 \\
\end{longtable}
\endgroup

\vspace{0.3cm}

\begingroup
\small
\renewcommand{\arraystretch}{1.25}
\begin{longtable}{@{} p{3.5cm} p{10.9cm} @{}}
\caption{{Equipo de cómputo del tercer desarrollador}} \label{tab:equipo_dev3} \\
\toprule
\textbf{Nombre} & \textbf{Descripción} \\
\midrule
\endfirsthead

\multicolumn{2}{@{}l}{\textbf{Tabla \ref{tab:equipo_dev3}} \textit{(Continuación)}} \\[0.4em]
\toprule
\textbf{Nombre} & \textbf{Descripción} \\
\midrule
\endhead

\midrule
\multicolumn{2}{r}{\textit{Continúa en la siguiente página\dots}} \\
\endfoot

\bottomrule
\multicolumn{2}{@{}p{\textwidth}}{\small \textit{Nota.} Computadora utilizada como equipo principal de trabajo por uno de los desarrolladores.} \\
\endlastfoot

Característica de la PC: & AMD Ryzen 7 5700G - 20.0 GB RAM - 424 GB Almacenamiento \\
Valor inicial: & S/ 2,600.00 \\
Vida Util en años: & 5 años \\
Valor residual: & S/ 520 \\
\end{longtable}
\endgroup

Para estimar el costo asociado al uso de las computadoras se aplica el método de depreciación lineal. Este método distribuye el valor depreciable del equipo en partes iguales a lo largo de su vida útil.

\begin{equation}
\text{Depreciación} =
\frac{\text{Valor inicial} - \text{Valor residual}}{\text{Vida útil en años}}
\end{equation}

\vspace{0.3cm}

\begingroup
\small
\renewcommand{\arraystretch}{1.25}
\begin{longtable}{@{} p{4.5cm} p{2.2cm} r c r c l @{}}
\caption{{Depreciación de los equipos de cómputo de cada desarrollador}} \label{tab:depreciacion_equipos} \\
\toprule
\textbf{Nombre} & \multicolumn{6}{l}{\textbf{Cálculo}} \\
\midrule
\endfirsthead

\multicolumn{7}{@{}l}{\textbf{Tabla \ref{tab:depreciacion_equipos}} \textit{(Continuación)}} \\[0.4em]
\toprule
\textbf{Nombre} & \multicolumn{6}{l}{\textbf{Cálculo}} \\
\midrule
\endhead

\midrule
\multicolumn{7}{r}{\textit{Continúa en la siguiente página\dots}} \\
\endfoot

\bottomrule
\multicolumn{7}{@{}p{\textwidth}}{\small \textit{Nota.} Cálculo realizado mediante el método de depreciación lineal, considerando el valor depreciable (Valor inicial menos valor residual) entre los 5 años de vida útil.} \\
\endlastfoot

Computadora del desarrollador 1 & Depreciación: & S/ 448.00 & = & S/ 2,240.00 & / & 5 años \\
Computadora del desarrollador 2 & Depreciación: & S/ 256.00 & = & S/ 1,280.00 & / & 5 años \\
Computadora del desarrollador 3 & Depreciación: & S/ 416.00 & = & S/ 2,080.00 & / & 5 años \\
\end{longtable}
\endgroup

De acuerdo con los cálculos anteriores, la depreciación acumulada de las tres computadoras principales asciende a \textbf{S/ 280.00} durante los tres meses de desarrollo (prorrateado a S/ 93.33 mensuales o S/ 1,120.00 anuales en conjunto).

\vspace{0.3cm}

\begingroup
\small
\renewcommand{\arraystretch}{1.25}
\begin{longtable}{@{} p{3.5cm} r c r c r @{}}
\caption{{Depreciación anual del equipo de cómputo del Desarrollador 1}} \label{tab:deprec_dev1_anios} \\
\toprule
\textbf{Valor} & \textbf{Valor residual} & & \textbf{Depreciación} & & \textbf{Resultado} \\
\midrule
\endfirsthead

\multicolumn{6}{@{}l}{\textbf{Tabla \ref{tab:deprec_dev1_anios}} \textit{(Continuación)}} \\[0.4em]
\toprule
\textbf{Valor} & \textbf{Valor residual} & & \textbf{Depreciación} & & \textbf{Resultado} \\
\midrule
\endhead

\midrule
\multicolumn{6}{r}{\textit{Continúa en la siguiente página\dots}} \\
\endfoot

\bottomrule
\multicolumn{6}{@{}p{\textwidth}}{\small \textit{Nota.} Muestra el proceso de depreciación lineal anual para el primer equipo hasta alcanzar su valor residual.} \\
\endlastfoot

1 año Valor: & S/ 2,240.00 & & & & \\
2 año Valor: & S/ 2,240.00 & - & S/ 448.00 & = & S/ 1,792.00 \\
3 año Valor: & S/ 1,792.00 & - & S/ 448.00 & = & S/ 1,344.00 \\
4 año Valor: & S/ 1,344.00 & - & S/ 448.00 & = & S/ 896.00 \\
5 año Valor: & S/ 896.00 & - & S/ 448.00 & = & S/ 448.00 \\
- & S/ 448.00 & - & S/ 448.00 & = & S/ 0.00 \\
\end{longtable}
\endgroup

\vspace{0.3cm}

\begingroup
\small
\renewcommand{\arraystretch}{1.25}
\begin{longtable}{@{} p{3.5cm} r c r c r @{}}
\caption{{Depreciación anual del equipo de cómputo del Desarrollador 2}} \label{tab:deprec_dev2_anios} \\
\toprule
\textbf{Valor} & \textbf{Valor residual} & & \textbf{Depreciación} & & \textbf{Resultado} \\
\midrule
\endfirsthead

\multicolumn{6}{@{}l}{\textbf{Tabla \ref{tab:deprec_dev2_anios}} \textit{(Continuación)}} \\[0.4em]
\toprule
\textbf{Valor} & \textbf{Valor residual} & & \textbf{Depreciación} & & \textbf{Resultado} \\
\midrule
\endhead

\midrule
\multicolumn{6}{r}{\textit{Continúa en la siguiente página\dots}} \\
\endfoot

\bottomrule
\multicolumn{6}{@{}p{\textwidth}}{\small \textit{Nota.} Muestra el proceso de depreciación lineal anual para el segundo equipo hasta alcanzar su valor residual.} \\
\endlastfoot

1 año Valor: & S/ 1,280.00 & & & & \\
2 año Valor: & S/ 1,280.00 & - & S/ 256.00 & = & S/ 1,024.00 \\
3 año Valor: & S/ 1,024.00 & - & S/ 256.00 & = & S/ 768.00 \\
4 año Valor: & S/ 768.00 & - & S/ 256.00 & = & S/ 512.00 \\
5 año Valor: & S/ 512.00 & - & S/ 256.00 & = & S/ 256.00 \\
- & S/ 256.00 & - & S/ 256.00 & = & S/ 0.00 \\
\end{longtable}
\endgroup

\vspace{0.3cm}

\begingroup
\small
\renewcommand{\arraystretch}{1.25}
\begin{longtable}{@{} p{3.5cm} r c r c r @{}}
\caption{{Depreciación anual del equipo de cómputo del Desarrollador 3}} \label{tab:deprec_dev3_anios} \\
\toprule
\textbf{Valor} & \textbf{Valor residual} & & \textbf{Depreciación} & & \textbf{Resultado} \\
\midrule
\endfirsthead

\multicolumn{6}{@{}l}{\textbf{Tabla \ref{tab:deprec_dev3_anios}} \textit{(Continuación)}} \\[0.4em]
\toprule
\textbf{Valor} & \textbf{Valor residual} & & \textbf{Depreciación} & & \textbf{Resultado} \\
\midrule
\endhead

\midrule
\multicolumn{6}{r}{\textit{Continúa en la siguiente página\dots}} \\
\endfoot

\bottomrule
\multicolumn{6}{@{}p{\textwidth}}{\small \textit{Nota.} Muestra el proceso de depreciación lineal anual para el tercer equipo hasta alcanzar su valor residual.} \\
\endlastfoot

1 año Valor: & S/ 2,080.00 & & & & \\
2 año Valor: & S/ 2,080.00 & - & S/ 416.00 & = & S/ 1,664.00 \\
3 año Valor: & S/ 1,664.00 & - & S/ 416.00 & = & S/ 1,248.00 \\
4 año Valor: & S/ 1,248.00 & - & S/ 416.00 & = & S/ 832.00 \\
5 año Valor: & S/ 832.00 & - & S/ 416.00 & = & S/ 416.00 \\
- & S/ 416.00 & - & S/ 416.00 & = & S/ 0.00 \\
\end{longtable}
\endgroup

\subsection{El Costo de Licencia de Software y Depreciación}

Para el adecuado funcionamiento del equipo de desarrollo, se adquirieron licenciamientos de software y recursos de equipamiento complementario. Además de las computadoras principales utilizadas por los desarrolladores, se consideran herramientas esenciales y de nivel profesional para la creación, pruebas integrales, despliegue continuo y mantenimiento del sistema HelpDesk. El uso de licencias oficiales y frameworks de desarrollo garantiza entornos de trabajo estables, seguros y compatibles con los estándares de la industria de software, permitiendo una integración continua y despliegues confiables en el entorno de producción. Asimismo, se contempla la depreciación proporcional de computadoras adicionales de soporte empleadas durante el ciclo de innovación, las cuales complementan el parque informático principal y aseguran que el equipo de trabajo cuente con recursos hardware suficientes para ejecutar los procesos de desarrollo, pruebas de carga y validación de funcionalidades sin interferir en el uso cotidiano de los equipos principales. Estos recursos de licenciamiento y equipamiento son indispensables para mantener la continuidad del equipo de trabajo, garantizar la integración de los frameworks de trabajo colaborativo y cumplir con los plazos del proyecto de innovación, ya que la ausencia de herramientas profesionales y equipos de soporte adecuados implicaría retrasos significativos en los tiempos de desarrollo, pruebas y validación del sistema. La \textbf{Tabla 30} detalla los importes correspondientes a dichos elementos, considerando licencias de herramientas de desarrollo y la depreciación proporcional de computadoras adicionales de soporte empleadas durante el ciclo de innovación, las cuales son indispensables para mantener la continuidad del equipo de trabajo y la integración de los frameworks de trabajo colaborativo.

\begingroup
\small
\renewcommand{\arraystretch}{1.25}
\begin{longtable}{@{} p{5.0cm} p{2.0cm} p{3.2cm} p{2.5cm} p{1.7cm} @{}}
\caption{{Costo de licenciamiento de software y equipamiento adicional}} \label{tab:licenciamiento_adicional} \\
\toprule
\textbf{Descripción} & \textbf{Cantidad} & \textbf{Costo unitario (S/.)} & \textbf{Total (S/.)} & \textbf{Dólar \$} \\
\midrule
\endfirsthead

\multicolumn{5}{@{}l}{\textbf{Tabla \ref{tab:licenciamiento_adicional}} \textit{(Continuación)}} \\[0.4em]
\toprule
\textbf{Descripción} & \textbf{Cantidad} & \textbf{Costo unitario (S/.)} & \textbf{Total (S/.)} & \textbf{Dólar \$} \\
\midrule
\endhead

\midrule
\multicolumn{5}{r}{\textit{Continúa en la siguiente página\dots}} \\
\endfoot

\bottomrule
\multicolumn{5}{@{}p{\textwidth}}{\small \textit{Nota.} Importes de licenciamiento y equipamiento complementario. Conversión a dólares utilizando un tipo de cambio referencial de S/ 3.42.} \\
\endlastfoot
Visual Studio Code & 7 & S/ 00.00 & S/ 00.00 & \$ 00.00 \\
Docker Compose & 2 & S/ 00.00 & S/ 00.00 & \$ 00.00 \\
Frameworks(NestJs, Vite 8.2 + React) & 7 & S/ 00.00 & S/ 00.00 & \$ 00.00 \\
Depreciación PC adicional 1 & 3 & S/ 395.83 & S/ 1,187.49 & \$ 342.98 \\
Depreciación PC adicional 2 & 2 & S/ 450.00 & S/ 900.00 & \$ 259.95 \\
Depreciación PC adicional 3 & 2 & S/ 410.00 & S/ 820.00 & \$ 236.84 \\
\textbf{Total} & & & \textbf{S/ 2,907.49} & \textbf{\$ 839.77} \\
\end{longtable}
\endgroup

\section{Otros costos de Implementación de mejora}
\subsection{El Costo de Servicio de consumo de Luz}

Para estimar el consumo eléctrico se considera un promedio de 230 kWh mensuales y una tarifa de S/ 0.65 por kWh.

\begin{equation}
\text{Costo mensual} = 230\text{ kWh} \times \text{S/ 0.65 por kWh} = \text{S/ 149.50}
\end{equation}

\vspace{0.3cm}

\begingroup
\small
\renewcommand{\arraystretch}{1.25}
\begin{longtable}{@{} p{5.5cm} p{2.5cm} p{2.5cm} p{2.0cm} p{1.9cm} @{}}
\caption{{Cálculo del consumo de energía eléctrica durante el proyecto}} \label{tab:consumo_electrico} \\
\toprule
\textbf{Concepto} & \textbf{Consumo} & \textbf{Tarifa} & \textbf{Duración} & \textbf{Total (S/.)} \\
\midrule
\endfirsthead

\multicolumn{5}{@{}l}{\textbf{Tabla \ref{tab:consumo_electrico}} \textit{(Continuación)}} \\[0.4em]
\toprule
\textbf{Concepto} & \textbf{Consumo} & \textbf{Tarifa} & \textbf{Duración} & \textbf{Total (S/.)} \\
\midrule
\endhead

\midrule
\multicolumn{5}{r}{\textit{Continúa en la siguiente página\dots}} \\
\endfoot

\bottomrule
\multicolumn{5}{@{}p{\textwidth}}{\small \textit{Nota.} Estimación del consumo eléctrico asignado al período de desarrollo en J\&P Periféricos S.A.C.} \\
\endlastfoot

Consumo de energía eléctrica & 230 kWh/mes & S/ 0.65/kWh & 3 meses & 448.50 \\
\textbf{Costo total} & & & \textbf{3 meses} & \textbf{448.50} \\
\end{longtable}
\endgroup

\subsection{Lista de gastos adicionales}

Se incluye también el servicio de internet Movistar de 400 Mbps con un costo mensual de S/ 89.90 durante tres meses.

\vspace{0.3cm}

\begingroup
\small
\renewcommand{\arraystretch}{1.25}
\begin{longtable}{@{} p{5.0cm} p{1.8cm} p{3.2cm} p{2.5cm} p{1.9cm} @{}}
\caption{{Gastos de operación y servicios utilizados durante el proyecto}} \label{tab:gastos_operacion} \\
\toprule
\textbf{Descripción} & \textbf{Meses} & \textbf{Precio mensual (S/.)} & \textbf{Total (S/.)} & \textbf{Dólar \$} \\
\midrule
\endfirsthead

\multicolumn{5}{@{}l}{\textbf{Tabla \ref{tab:gastos_operacion}} \textit{(Continuación)}} \\[0.4em]
\toprule
\textbf{Descripción} & \textbf{Meses} & \textbf{Precio mensual (S/.)} & \textbf{Total (S/.)} & \textbf{Dólar \$} \\
\midrule
\endhead

\midrule
\multicolumn{5}{r}{\textit{Continúa en la siguiente página\dots}} \\
\endfoot

\bottomrule
\multicolumn{5}{@{}p{\textwidth}}{\small \textit{Nota.} Gastos operativos de energía e internet por tres meses. Conversión a dólares utilizando el tipo de cambio referencial de S/ 3.42.} \\
\endlastfoot

Servicio de energía eléctrica & 3 & 149.50 & 448.50 & \$ 131.14 \\
Servicio de internet (Movistar 400 Mbps) & 3 & 89.90 & 269.70 & \$ 78.86 \\
\textbf{Total} & & & \textbf{718.20} & \textbf{\$ 210.00} \\
\end{longtable}
\endgroup

\section{Costo total de la Implementación de Mejora / Innovación / Creatividad}

A continuación se consolidan todos los rubros en el presupuesto general de la implementación del sistema.

\vspace{0.3cm}

\begingroup
\small
\renewcommand{\arraystretch}{1.25}
\begin{longtable}{@{} p{9.4cm} p{2.5cm} p{2.5cm} @{}}
\caption{{Presupuesto general de implementación del sistema HelpDesk}} \label{tab:presupuesto_general} \\
\toprule
\textbf{Descripción} & \textbf{Total (S/.)} & \textbf{Dólar \$} \\
\midrule
\endfirsthead

\multicolumn{3}{@{}l}{\textbf{Tabla \ref{tab:presupuesto_general}} \textit{(Continuación)}} \\[0.4em]
\toprule
\textbf{Descripción} & \textbf{Total (S/.)} & \textbf{Dólar \$} \\
\midrule
\endhead

\midrule
\multicolumn{3}{r}{\textit{Continúa en la siguiente página\dots}} \\
\endfoot

\bottomrule
\multicolumn{3}{@{}p{\textwidth}}{\small \textit{Nota.} Presupuesto consolidado para la implementación del sistema utilizando un tipo de cambio referencial de S/ 3.42.} \\
\endlastfoot

Costo de materiales y servicios digitales & 543.00 & 158.77 \\
Costo de mano de obra & 13,410.00 & 3,921.05 \\
Costo de máquinas, herramientas y equipos & 3,582.99 & 1,047.66 \\
Otros costos de operación y conectividad & 718.20 & 210.00 \\
\textbf{Total general de implementación} & \textbf{18,254.19} & \textbf{5,337.48} \\
\end{longtable}
\endgroup

\vspace{0.2cm}

La inversión total de \textbf{S/ 18,254.19} (US\$ 5,337.48) en la implementación del sistema \textit{HelpDesk} refleja una estructura financiera concentrada en un 73.5\% en mano de obra especializada, un 19.6\% en máquinas, herramientas y equipos, un 3.9\% en otros costos de operación y conectividad, y un 3.0\% en materiales y servicios digitales. Cada rubro ha sido considerado bajo criterios de eficiencia presupuestaria, aprovechamiento de recursos internos y operatividad tecnológica para J\&P Periféricos S.A.C.